\clearpage
\section{The surviving leading symplectic layer}\label{app:odd_capacity}

Status: ordinary proof from the retained Bushnell--Kutzko theorems and,
for the source-defined candidate, the substantial published
Bushnell--Henniart arithmetic induction/descent constructor. No Lean
build or bounded computation is used. The completed independent review
is retained in the accompanying audit files.

\hypertarget{odd_capacity-the-useful-simplification}{%
\subsection{The useful simplification}\label{odd_capacity-the-useful-simplification}}

Full critical-profile rigidity is unnecessary for the trivial-stabilizer
finite attachment contract. An even leading level produces a symplectic
subspace that survives every later simple approximation. For the
independent rank-p Heisenberg tensor sources, this single subspace
already has enough dimension to hold every source plane, however many
there are.

The argument below concerns the actual type in the source-defined
arithmetic candidate. It does not identify that type with a separately
guessed binder parameter.

\hypertarget{odd_capacity-leading-layer-theorem}{%
\subsection{1. Leading-layer theorem}\label{odd_capacity-leading-layer-theorem}}

Let F have residue field k of cardinality q and residue characteristic
p.~Let pi be an actual irreducible supercuspidal representation of
GL\_n(F) with torsion number one, where n=p\^{}r. Suppose

\begin{Verbatim}[fontsize=\small,breaklines=true,breakanywhere=true]
a(pi)=n+N,     N>0 even,     v_p(N)=r−1.
\end{Verbatim}

Then an actual maximal simple type in pi, with Heisenberg quotient
W=J\textsuperscript{1/H}1, contains a nondegenerate symplectic
F\_p-subspace S of dimension

\begin{Verbatim}[fontsize=\small,breaklines=true,breakanywhere=true]
[k:F_p] (p−1)p^(r−1).
\end{Verbatim}

For k=F\_p this is exactly (p−1)p\^{}(r−1). The subspace is the image of
the concrete congruence group U\^{}(N/2)(A) in J\textsuperscript{1/H}1.
It is present regardless of the later critical exponents, their
parities, or the number of intermediate approximation fields.

\hypertarget{odd_capacity-actual-parameter-field-and-leading-level}{%
\subsubsection{1.1. Actual parameter field and leading
level}\label{odd_capacity-actual-parameter-field-and-leading-level}}

By BK maximal-type existence, pi contains a maximal simple type. Its
positive level follows from a(pi)\textgreater n.~Write its stratum as
{[}A,N\_A,0,beta{]}, with parameter field K=F{[}beta{]}. BK 6.2.5 gives

\begin{Verbatim}[fontsize=\small,breaklines=true,breakanywhere=true]
t(pi)=n/e(K/F).
\end{Verbatim}

Thus e(K/F)=n.~Since {[}K:F{]} divides n, K/F is totally ramified of
degree n and has residue degree one. The principal-order assertion after
BK 6.2.5 gives e(A)=n.~Accordingly A is the complete principal
hereditary order. The established conductor formula

\begin{Verbatim}[fontsize=\small,breaklines=true,breakanywhere=true]
a(pi)=n+n N_A/e(A)
\end{Verbatim}

now implies N\_A=N. This is the same source-checked formula used in
Block165, originating in BH1999 Section 6.1 Lemma 2 and retained in
\path{PRINCIPAL_ORDER_CONDUCTOR161.md}.

Put P=rad(A), U\textsuperscript{a(A)=1+P}a for a\textgreater=1, and
m=N/2.

\hypertarget{odd_capacity-the-terminal-minimal-approximation-has-degree-p}{%
\subsubsection{1.2. The terminal minimal approximation has degree
p}\label{odd_capacity-the-terminal-minimal-approximation-has-degree-p}}

Apply the defining-sequence construction BK 2.4.1--2 to beta. Write
gamma for its terminal minimal approximation. Every transition has a
critical level c\textless N, and gamma has the same A-level N as beta.
All approximation residue degrees divide f(K/F)=1, and all approximation
ramification degrees divide n.~Therefore E=F{[}gamma{]} is totally
ramified of degree d dividing n.

The scalar case is impossible: a scalar of F has A-valuation divisible
by n, whereas v\_p(N)=r−1. Thus gamma is noncentral and minimal in the
ordinary field sense. Its field valuation is

\begin{Verbatim}[fontsize=\small,breaklines=true,breakanywhere=true]
−v_E(gamma)=Nd/n.
\end{Verbatim}

Minimality, by BK 1.4.14--15, requires gcd(Nd/n,d)=1. Since d is a
p-power and v\_p(N)=r−1, integrality first gives p\textbar d, while
coprimality forces d=p.~Hence E/F is totally ramified of degree p.

Set B=End\_E(F\^{}n), B\_order=A∩B, and Q=rad(B\_order). The same
complete lattice chain has period n/p over E. Each P-grade has
k-dimension n and each Q-grade has k-dimension n/p.~In particular,

\begin{Verbatim}[fontsize=\small,breaklines=true,breakanywhere=true]
dim_k P^m/(Q^m+P^(m+1))=n−n/p.
\end{Verbatim}

No assertion that E is literally a subfield of the final field K is
needed. It is the actual minimal approximation field in the defining
sequence.

\hypertarget{odd_capacity-the-high-subgroup-survives-every-approximation}{%
\subsubsection{1.3. The high subgroup survives every
approximation}\label{odd_capacity-the-high-subgroup-survives-every-approximation}}

The minimal definitions BK 3.1.7--8 give

\begin{Verbatim}[fontsize=\small,breaklines=true,breakanywhere=true]
J_order(gamma)=B_order+P^m,
H_order(gamma)=B_order+P^(m+1).
\end{Verbatim}

At a nonminimal transition beta\_i→beta\_(i+1), with critical level
c\_i\textless N, the recursive J-order definition uses the threshold
ceil(c\_i/2). Since ceil(c\_i/2)\textless=m, the inclusion
P\^{}m⊂J\_order(beta\_(i+1)) implies P\^{}m⊂J\_order(beta\_i). Starting
at gamma, induction proves

\begin{Verbatim}[fontsize=\small,breaklines=true,breakanywhere=true]
U^m(A)⊂J^1(beta).
\end{Verbatim}

For H, BK 3.1.9(ii) states that
H\textsuperscript{a(beta\_i)=H}a(beta\_(i+1)) as soon as

\begin{Verbatim}[fontsize=\small,breaklines=true,breakanywhere=true]
a>=floor(c_i/2)+1.
\end{Verbatim}

Every c\_i is at most 2m−1, so this threshold is at most m. Consequently

\begin{Verbatim}[fontsize=\small,breaklines=true,breakanywhere=true]
H^m(beta)=H^m(gamma).
\end{Verbatim}

Here H\^{}m means H\textsuperscript{1∩U}m(A). Therefore the image S of
U\^{}m(A) in W is canonically

\begin{Verbatim}[fontsize=\small,breaklines=true,breakanywhere=true]
S = U^m(A)/H^m(beta)
  = U^m(A)/H^m(gamma)
  ≅ P^m/(Q^m+P^(m+1)).
\end{Verbatim}

The last group identification uses m\textgreater=1, so multiplication is
additive modulo P\^{}(m+1). Thus \textbar S\textbar=q\^{}(n−n/p).

\hypertarget{odd_capacity-its-actual-commutator-form-is-nondegenerate}{%
\subsubsection{1.4. Its actual commutator form is
nondegenerate}\label{odd_capacity-its-actual-commutator-form-is-nondegenerate}}

Choose the actual simple character theta contained in the type of pi. By
the standard deep-restriction property of simple characters, theta
restricts to psi\_beta on U\^{}(m+1)(A). This follows directly from BK
3.2.1--4 by induction along the defining sequence.

For u,v in U\^{}m(A), their group commutator is in
U\textsuperscript{(2m)(A)=U}N(A), which lies in U\^{}(m+1)(A). The
terminal approximation satisfies

\begin{Verbatim}[fontsize=\small,breaklines=true,breakanywhere=true]
beta−gamma∈P^(1−N),
\end{Verbatim}

because every critical level in the sequence is less than N. Hence

\begin{Verbatim}[fontsize=\small,breaklines=true,breakanywhere=true]
(beta−gamma)([u,v]−1)∈P.
\end{Verbatim}

With the standard level-one additive character psi\_A=psi\_F∘tr\_A/F,
psi\_A is trivial on P. We therefore have the exact equality

\begin{Verbatim}[fontsize=\small,breaklines=true,breakanywhere=true]
theta([u,v]) = psi_beta([u,v]) = psi_gamma([u,v]).
\end{Verbatim}

The right side is the Heisenberg commutator pairing for the minimal
simple stratum of gamma. Its quotient
J\textsuperscript{1(gamma)/H}1(gamma) is precisely
U\textsuperscript{m(A)/H}m(gamma), because
J\textsuperscript{1(gamma)=U}1(B\_order)U\^{}m(A), with
U\textsuperscript{1(B\_order)⊂H}1(gamma). BK's nondegeneracy theorem for
simple-character commutators therefore says that the form on S is
nondegenerate.

For general k, this is a nondegenerate alternating F\_p-form of
dimension \href{n−n/p}{k:F\_p}. For k=F\_p, it has dimension n−n/p.~This
proves the theorem.

The argument uses no self-pair conductor and no determination of the
later critical levels. Those invariants remain useful for finer
arithmetic identification, but they are not needed for this
finite-sector capacity claim.

\hypertarget{odd_capacity-arbitrary-length-odd-prime-source-family}{%
\subsection{2. Arbitrary-length odd-prime source
family}\label{odd_capacity-arbitrary-length-odd-prime-source-family}}

Fix an odd prime p and F=F\_p((t)). Let the complete marked Weil source
rho be an independent tensor product of r rank-p wild Heisenberg factors
at distinct positive levels

\begin{Verbatim}[fontsize=\small,breaklines=true,breakanywhere=true]
0<h_1<...<h_r,     p does not divide h_i.
\end{Verbatim}

Here independence means the stipulated balanced-central-product
description of the wild image is valid, and rho\textbar P\_F is
irreducible. We do not infer independence from the inequalities alone.
Suppose the retained Tsushima source calculation supplies the usual
highest slope

\begin{Verbatim}[fontsize=\small,breaklines=true,breakanywhere=true]
slope(rho)=(p+1)h_r/p.
\end{Verbatim}

Then rank(rho)=n=p\^{}r and

\begin{Verbatim}[fontsize=\small,breaklines=true,breakanywhere=true]
Swan(rho)=N=(p+1)p^(r−1)h_r.
\end{Verbatim}

Let Pi\_BH(rho) be the published source-defined arithmetic
induction/descent candidate, with its full prescribed central character,
as audited in Block165. The construction is applied to the complete
actual rho, including its scalar normalization. Its codomain is the
torsion-one supercuspidal class, and its standard epsilon-function
property gives a(Pi\_BH(rho))=n+N. These are uses of the substantial
classical constructor, not consequences of finite Heisenberg algebra
alone and not a selection of an inverse LLC image.

For odd p, N is even and v\_p(N)=r−1. The leading-layer theorem
therefore puts a symplectic space of dimension

\begin{Verbatim}[fontsize=\small,breaklines=true,breakanywhere=true]
D=(p−1)p^(r−1)
\end{Verbatim}

inside the actual Heisenberg quotient of a type in this same
Pi\_BH(rho).

The source packet is the orthogonal sum of r planes, so its symplectic
dimension is 2r. The elementary inequality

\begin{Verbatim}[fontsize=\small,breaklines=true,breakanywhere=true]
(p−1)p^(r−1)>=2r
\end{Verbatim}

holds for every odd p and r\textgreater=1: it suffices that
3\^{}(r−1)\textgreater=r, which follows by induction. Thus the actual
candidate's leading symplectic layer alone can receive the entire source
packet.

\hypertarget{odd_capacity-full-characters-and-recovery-for-declared-c1}{%
\subsubsection{2.1. Full characters and recovery for declared
C=\{1\}}\label{odd_capacity-full-characters-and-recovery-for-declared-c1}}

Work under the same finite scalar-pushout contract as
\path{SAME_CANDIDATE_ATTACHMENT165.md}: declared stabilizer C=\{1\}, one
vertex, and full source character extensions retained, with no
requirement of literal old matrix-group inclusion or equality of old
numerical depths with host depths.

Let P\_src be the balanced C×-central product of the r supplied full
source character extensions. Its quotient V\_src has dimension 2r over
F\_p and its central-character-one Heisenberg representation has
dimension p\^{}r=n.~Choose a symplectic embedding V\_src→S⊂W.

The C×-pushout of the actual type character, restricted to this
subspace, has the same alternating commutator as P\_src. The retained
extension theorem supplies a center-fixing isomorphism from P\_src to
that restriction. Consequently every source h∈H\_i\^{}1 is carried to
the scalar theta\_i(h), preserving the full character value. No shallow
source values are discarded.

Let V be the image and write W=V⊥V\_perp. Stone--von Neumann gives

\begin{Verbatim}[fontsize=\small,breaklines=true,breakanywhere=true]
eta_W ≅ eta_src⊗eta_perp.
\end{Verbatim}

Choose a line in eta\_perp and the corresponding rank-n projector e in
the actual finite Heisenberg image algebra. Then e eta\_W≅eta\_src.
Applying the retained Block159 induction/idempotent theorem to the
unchanged actual extended type Lambda in Pi\_BH(rho) recovers Lambda and
then this exact Pi\_BH(rho). All source normalization and the actual
noncompact scalar remain fixed.

This proves same-candidate finite attachment for arbitrary r in this
C=\{1\} odd-prime family. Original H5, when supplied for this same
candidate, remains exactly the original conditional comparison premise.
We do not prove it here.

\hypertarget{odd_capacity-first-new-rank-gl_27}{%
\subsubsection{2.2. First new rank: GL\_27}\label{odd_capacity-first-new-rank-gl_27}}

For p=3,r=3, the source has rank27 and V\_src has dimension6. Put h=h\_3
with 3∤h. The actual candidate has

\begin{Verbatim}[fontsize=\small,breaklines=true,breakanywhere=true]
N=36h,       m=18h,
dim_F3 S=27−9=18.
\end{Verbatim}

Thus its high layer contains the entire three-plane source packet. At
least a 3\^{}9-dimensional Heisenberg contribution is present, while the
required source sector has dimension27=3\^{}3. The six surplus
symplectic pairs can be projected onto a line in their representation,
together with any additional lower layers. This proves the required
capacity and sector realization for the new rank27 family without
solving its two possible lower critical levels.

In particular, for (h\_1,h\_2,h\_3)=(1,2,4), N=144 and the same high
subspace has dimension18. These numbers are a specialization of the
universal proof, not a numerical experiment.

\hypertarget{odd_capacity-why-trying-to-force-every-critical-level-was-unnecessary}{%
\subsection{3. Why trying to force every critical level was
unnecessary}\label{odd_capacity-why-trying-to-force-every-critical-level-was-unnecessary}}

At rank27, a hypothetical complete degree pattern 27→9→3 has leading
level N=36h\_3 and lower critical levels a\textless b\textless N. The
standard full-profile conductor count would give

\begin{Verbatim}[fontsize=\small,breaklines=true,breakanywhere=true]
Swan(End rho)=18N+6b+2a.
\end{Verbatim}

Comparing with the stipulated source formula 8(h\_1+9h\_2+81h\_3) only
gives

\begin{Verbatim}[fontsize=\small,breaklines=true,breakanywhere=true]
a+3b=4h_1+36h_2.
\end{Verbatim}

That single linear relation does not determine both a and b, and by
itself does not force each parity. Nor does it establish that all three
degrees occur rather than skipping degree9. This observation is only a
warning about the insufficient equations; it is not presented as an
arithmetic counterexample or as existence of every numerical profile
satisfying the equation.

The leading-layer theorem bypasses these unresolved profile choices. It
supplies a fixed nondegenerate high subspace in every actual candidate
with the two required invariants.

\hypertarget{odd_capacity-dyadic-capacity-boundary}{%
\subsection{4. Dyadic capacity
boundary}\label{odd_capacity-dyadic-capacity-boundary}}

The leading-layer theorem itself is valid at p=2. If an actual
torsion-one candidate of rank2\^{}r has

\begin{Verbatim}[fontsize=\small,breaklines=true,breakanywhere=true]
N=3·2^(r−1)h,     h odd,     r>=2,
\end{Verbatim}

then its high subspace has dimension 2\^{}(r−1) over F\_2. This alone is
at least 2r for r\textgreater=4, with equality at r=4. It does not
supply enough capacity by itself for r=2 or3. For r=1, N is odd and this
theorem does not apply.

This is a capacity statement only. Applying it to a dyadic source
requires the actual source family, its arithmetic candidate and whatever
quadratic/phase markings the original contract requires. No universal
dyadic attachment theorem is claimed from dimension alone.

\hypertarget{odd_capacity-boundaries-that-remain}{%
\subsection{5. Boundaries that remain}\label{odd_capacity-boundaries-that-remain}}

\begin{enumerate}
\def\labelenumi{\arabic{enumi}.}
\tightlist
\item
  The full character attachment here is through center-fixing
  scalar-pushout maps. Literal embeddings of prescribed old arithmetic
  matrix groups are not proved.
\item
  No nontrivial prescribed stabilizer action is attached. Capacity does
  not imply equivariant placement.
\item
  All planes may be placed in the leading host layer. Preservation of
  supplied numerical ramification depths, or a specified map into
  individual host depth layers, would be an additional requirement and
  is not claimed.
\item
  General packet transport graphs are not handled here. The statement is
  for C=\{1\} and one vertex.
\item
  The arithmetic induction/descent theorem is an explicit major external
  input. This is a new capacity/attachment consequence for its actual
  output, not a new elementary reconstruction of that classical
  constructor.
\item
  H5 remains conditional and refers to that unchanged actual candidate.
  No additional pair-factor premise is introduced.
\item
  No Lean verification is claimed.
\end{enumerate}

\hypertarget{odd_capacity-primary-locators-and-retained-dependencies}{%
\subsection{6. Primary locators and retained
dependencies}\label{odd_capacity-primary-locators-and-retained-dependencies}}

The authorized primary BK text was read at these exact results:
1.4.14--15 (minimality), 2.4.1--2 (defining sequence and degree
divisibility), 3.1.7--9 (recursive orders and high-depth equality),
3.2.1--4 (deep restriction of simple characters), 3.4.1 (nondegenerate
Heisenberg commutator), and 6.2.5 with the following paragraph (torsion
number and principal period).

The standard conductor and published arithmetic-constructor properties
are reused from the source-pinned Block165 audit and Block161 conductor
proof. The finite scalar-pushout and sector-recovery arguments are
exactly those retained and restated in
\path{SAME_CANDIDATE_ATTACHMENT165.md}. No entire copyrighted source is
copied into this continuation.
